\documentclass[11pt,a4paper]{article}
\usepackage[utf8]{inputenc}
\usepackage[T1]{fontenc}
\usepackage[spanish,es-minimal,es-noshorthands]{babel}
\usepackage{lmodern}
\usepackage[sfdefault,lining]{FiraSans}
\usepackage[a4paper,top=2.3cm,bottom=2.2cm,left=1.8cm,right=1.8cm,headheight=1.3cm,headsep=0.4cm,footskip=1cm]{geometry}
\usepackage{amsmath,amssymb}
\usepackage{icomma}
\usepackage{graphicx}
\usepackage[table]{xcolor}
\usepackage{tikz}
\usepackage{fontawesome5}
\usepackage{ccicons}
\usepackage{fancyhdr}
\usepackage{lastpage}
\usepackage{tasks}
\usepackage{needspace}
\usepackage[most]{tcolorbox}
\usepackage{array}
\usepackage[hidelinks]{hyperref}
\AtBeginDocument{\hypersetup{pdftitle={Ficha de ejercicios de la unidad 3: Ecuaciones e inecuaciones. 4º ESO},pdfsubject={Matemáticas B 4º ESO},pdfkeywords={ecuaciones de segundo grado, bicuadradas, Ruffini, factorizadas, fracciones algebraicas, radicales, exponenciales, logarítmicas, inecuaciones}}}

\definecolor{navy}{HTML}{1B2748}
\definecolor{accent}{HTML}{E8603A}
\definecolor{teal}{HTML}{1F9E90}
\definecolor{soft}{HTML}{EEF1F8}
\definecolor{muted}{HTML}{5D6685}

\setlength{\parindent}{0pt}
\setlength{\parskip}{3pt}
\settasks{label=\alph*), label-format=\bfseries\color{navy}, label-width=1.3em, item-indent=1.8em, column-sep=0.6em, after-item-skip=5pt, before-skip=3pt, after-skip=4pt}

% ---------- capçalera i peu ----------
\newcommand{\llicencia}{\href{https://creativecommons.org/licenses/by-nc-sa/4.0/deed.es}{\ccbyncsa}\ \ Licencia Creative Commons Atribución-NoComercial-CompartirIgual 4.0 Internacional (CC BY-NC-SA 4.0)}
\pagestyle{fancy}
\fancyhf{}
\fancyhead[L]{\raisebox{-0.35\height}{\includegraphics[height=0.95cm]{logo-gva.png}}}
\fancyhead[C]{\small\color{muted}Ficha de ejercicios · Unidad 3: Ecuaciones e inecuaciones · 4\textsuperscript{o} ESO B}
\fancyhead[R]{\raisebox{-0.35\height}{\includegraphics[height=0.95cm]{logo-vermellar.png}}}
\fancyfoot[L]{\scriptsize\color{muted}\llicencia}
\fancyfoot[R]{\small\color{muted}\thepage}
\renewcommand{\headrulewidth}{0.4pt}
\renewcommand{\headrule}{\color{muted!40}\hrule height \headrulewidth}
\fancypagestyle{primera}{\fancyhf{}\renewcommand{\headrulewidth}{0pt}
  \fancyfoot[L]{\scriptsize\color{muted}\llicencia}
  \fancyfoot[R]{\small\color{muted}\thepage}}

% ---------- nivells i exercicis ----------
\newcommand{\estrelles}[1]{{\color{accent}\ifcase#1\or\faStar\faStar[regular]\faStar[regular]\or\faStar\faStar\faStar[regular]\or\faStar\faStar\faStar\fi}}
\newcounter{exer}
\newcommand{\ex}[1]{\par\needspace{4\baselineskip}\medskip\refstepcounter{exer}\noindent{\color{navy}\bfseries\theexer.}\ \ {\small\estrelles{#1}}\ \ }

\newcommand{\seccio}[2]{\par\needspace{16\baselineskip}\bigskip
  \begin{tcolorbox}[colback=navy,colframe=navy,coltext=white,arc=2mm,boxrule=0pt,left=3mm,right=3mm,top=1.2mm,bottom=1.2mm]
  {\large\bfseries #1\quad #2}
  \end{tcolorbox}\vspace{-2pt}}

\newcommand{\exf}[1]{\par\medskip\refstepcounter{exer}\noindent{\color{navy}\bfseries\theexer.}\ \ {\small\estrelles{#1}}\ \ }
\newcolumntype{C}[1]{>{\centering\arraybackslash}p{#1}}
% quadre de recordatori (només a la fitxa d'exercicis)
\newcommand{\recorda}[1]{\par\smallskip\begin{tcolorbox}[colback=accent!7,colframe=accent!55,arc=2mm,boxrule=0.5pt,left=2.5mm,right=2.5mm,top=1mm,bottom=1mm]\small{\color{accent}\faLightbulb[regular]}\ \textbf{Recuerda:}\ #1\end{tcolorbox}}
\newcommand{\res}[1]{#1}\newcommand{\rest}[1]{#1}

\def\fmneed{11}% que el títol d'una secció no quede sol al final del full
% =====================================================================
% Capçalera FentMates de les fitxes: un color per a cada curs.
% Els colors són els de la pàgina inicial (LEVELS a assets/js/materials.js);
% el logo ve de _eines/logo/pdf/ (python3 _eines/logo/genera.py).
%
% Ús, al final del preàmbul de la unitat:
%   % =====================================================================
% Capçalera FentMates de les fitxes: un color per a cada curs.
% Els colors són els de la pàgina inicial (LEVELS a assets/js/materials.js);
% el logo ve de _eines/logo/pdf/ (python3 _eines/logo/genera.py).
%
% Ús, al final del preàmbul de la unitat:
%   % =====================================================================
% Capçalera FentMates de les fitxes: un color per a cada curs.
% Els colors són els de la pàgina inicial (LEVELS a assets/js/materials.js);
% el logo ve de _eines/logo/pdf/ (python3 _eines/logo/genera.py).
%
% Ús, al final del preàmbul de la unitat:
%   \input{../comu/fentmates.tex}
%   \fmcurs{1eso}
% i, al començament del document:
%   \capcalera{Matemàtiques}{1r ESO}{Fitxa d'exercicis · Unitat 2}{Divisibilitat}
% Opcional abans de \input: \def\fmneed{14} (espai mínim abans de cada secció).
% =====================================================================
\makeatletter
\def\fm@logodir{../../logo/pdf}
\providecommand{\fmneed}{10}
\def\fm@cursnom{}

% clau del curs, color i color del text damunt del color
\newcommand{\fm@curs}[3]{\@namedef{fm@color@#1}{#2}\@namedef{fm@text@#1}{#3}}
\fm@curs{1eso}{E6B84A}{navy}
\fm@curs{2eso}{3B82F6}{white}
\fm@curs{3eso}{FF7A50}{white}
\fm@curs{4eso-a}{A78BFA}{white}
\fm@curs{4eso-b}{6D28D9}{white}
\fm@curs{1batx-mi}{9B1C1C}{white}
\fm@curs{1batx-ccss}{EF6B6B}{white}
\fm@curs{2batx-mii}{1E6B47}{white}
\fm@curs{2batx-ccss}{2BB3A3}{white}

\newcommand{\fmcurs}[1]{%
  \@ifundefined{fm@color@#1}{\PackageError{fentmates}{Curs desconegut: #1}{}}{}%
  \edef\fm@tmp{\noexpand\definecolor{curs}{HTML}{\@nameuse{fm@color@#1}}%
               \noexpand\colorlet{curstext}{\@nameuse{fm@text@#1}}}\fm@tmp
  \colorlet{cursfosc}{curs!72!black}%
  \def\fm@logo{\fm@logodir/vertical-#1.pdf}%
  \def\fm@marca{\fm@logodir/marca-#1.pdf}}

% ---------- capçalera de la primera pàgina ----------
% #1 matèria (petit), #2 curs, #3 tipus de document i unitat, #4 títol
\newcommand{\capcalera}[4]{%
  \gdef\fm@cursnom{#2}%
  \par\noindent
  \begin{tcolorbox}[colback=curs,colframe=curs,boxrule=0pt,arc=2.2mm,width=3.5cm,height=2.15cm,
      valign=center,left=3mm,right=2mm,top=1mm,bottom=1mm,nobeforeafter,box align=center]
    \color{curstext}{\footnotesize #1}\par\vspace{1pt}{\fontsize{19}{21}\selectfont\bfseries #2\par}%
  \end{tcolorbox}\hspace{5mm}%
  \begin{minipage}[c]{\dimexpr\linewidth-3.5cm-5mm-2.9cm\relax}
    {\color{cursfosc}\small\bfseries\MakeUppercase{#3}}\par\vspace{2pt}
    {\color{navy}\fontsize{20}{23}\selectfont\bfseries #4\par}
  \end{minipage}\hfill
  \begin{minipage}[c]{2.6cm}\raggedleft\includegraphics[width=2.6cm]{\fm@logo}\end{minipage}\par
  \vspace{2pt}}

% ---------- barres de secció del color del curs ----------
\renewcommand{\seccio}[2]{\par\needspace{\fmneed\baselineskip}\bigskip
  \begin{tcolorbox}[colback=curs,colframe=curs,coltext=curstext,arc=2mm,boxrule=0pt,left=3mm,right=3mm,top=1.2mm,bottom=1.2mm]
  {\large\bfseries #1\quad #2}
  \end{tcolorbox}\vspace{-2pt}}

% ---------- estrelles de dificultat del color del curs (a Batxillerat no se'n posen) ----------
\@ifundefined{estrelles}{}{%
  \renewcommand{\estrelles}[1]{{\color{curs}\ifcase#1\or\faStar\faStar[regular]\faStar[regular]\or\faStar\faStar\faStar[regular]\or\faStar\faStar\faStar\fi}}}

% ---------- capçalera i peu de les pàgines ----------
% primera pàgina: logos institucionals al peu
\fancypagestyle{primera}{\fancyhf{}\renewcommand{\headrulewidth}{0pt}
  \fancyfoot[L]{\raisebox{-0.32\height}{\includegraphics[height=0.78cm]{logo-gva.png}}\hspace{5mm}%
    \raisebox{-0.32\height}{\includegraphics[height=0.78cm]{logo-vermellar.png}}}
  \fancyfoot[R]{\scriptsize\color{muted}\href{https://creativecommons.org/licenses/by-nc-sa/4.0/deed.ca}{\ccbyncsa}\ \ CC BY-NC-SA 4.0 · fentmates.org\qquad\small\thepage}}
% la resta: el símbol a l'esquerra i el curs a la dreta, amb el títol al centre com abans
\fancyhead[L]{\raisebox{-0.3\height}{\includegraphics[height=0.62cm]{\fm@marca}}}
\fancyhead[R]{\small\bfseries\color{cursfosc}\fm@cursnom}
\renewcommand{\headrule}{\color{curs!45}\hrule height \headrulewidth}
\makeatother

%   \fmcurs{1eso}
% i, al començament del document:
%   \capcalera{Matemàtiques}{1r ESO}{Fitxa d'exercicis · Unitat 2}{Divisibilitat}
% Opcional abans de \input: \def\fmneed{14} (espai mínim abans de cada secció).
% =====================================================================
\makeatletter
\def\fm@logodir{../../logo/pdf}
\providecommand{\fmneed}{10}
\def\fm@cursnom{}

% clau del curs, color i color del text damunt del color
\newcommand{\fm@curs}[3]{\@namedef{fm@color@#1}{#2}\@namedef{fm@text@#1}{#3}}
\fm@curs{1eso}{E6B84A}{navy}
\fm@curs{2eso}{3B82F6}{white}
\fm@curs{3eso}{FF7A50}{white}
\fm@curs{4eso-a}{A78BFA}{white}
\fm@curs{4eso-b}{6D28D9}{white}
\fm@curs{1batx-mi}{9B1C1C}{white}
\fm@curs{1batx-ccss}{EF6B6B}{white}
\fm@curs{2batx-mii}{1E6B47}{white}
\fm@curs{2batx-ccss}{2BB3A3}{white}

\newcommand{\fmcurs}[1]{%
  \@ifundefined{fm@color@#1}{\PackageError{fentmates}{Curs desconegut: #1}{}}{}%
  \edef\fm@tmp{\noexpand\definecolor{curs}{HTML}{\@nameuse{fm@color@#1}}%
               \noexpand\colorlet{curstext}{\@nameuse{fm@text@#1}}}\fm@tmp
  \colorlet{cursfosc}{curs!72!black}%
  \def\fm@logo{\fm@logodir/vertical-#1.pdf}%
  \def\fm@marca{\fm@logodir/marca-#1.pdf}}

% ---------- capçalera de la primera pàgina ----------
% #1 matèria (petit), #2 curs, #3 tipus de document i unitat, #4 títol
\newcommand{\capcalera}[4]{%
  \gdef\fm@cursnom{#2}%
  \par\noindent
  \begin{tcolorbox}[colback=curs,colframe=curs,boxrule=0pt,arc=2.2mm,width=3.5cm,height=2.15cm,
      valign=center,left=3mm,right=2mm,top=1mm,bottom=1mm,nobeforeafter,box align=center]
    \color{curstext}{\footnotesize #1}\par\vspace{1pt}{\fontsize{19}{21}\selectfont\bfseries #2\par}%
  \end{tcolorbox}\hspace{5mm}%
  \begin{minipage}[c]{\dimexpr\linewidth-3.5cm-5mm-2.9cm\relax}
    {\color{cursfosc}\small\bfseries\MakeUppercase{#3}}\par\vspace{2pt}
    {\color{navy}\fontsize{20}{23}\selectfont\bfseries #4\par}
  \end{minipage}\hfill
  \begin{minipage}[c]{2.6cm}\raggedleft\includegraphics[width=2.6cm]{\fm@logo}\end{minipage}\par
  \vspace{2pt}}

% ---------- barres de secció del color del curs ----------
\renewcommand{\seccio}[2]{\par\needspace{\fmneed\baselineskip}\bigskip
  \begin{tcolorbox}[colback=curs,colframe=curs,coltext=curstext,arc=2mm,boxrule=0pt,left=3mm,right=3mm,top=1.2mm,bottom=1.2mm]
  {\large\bfseries #1\quad #2}
  \end{tcolorbox}\vspace{-2pt}}

% ---------- estrelles de dificultat del color del curs (a Batxillerat no se'n posen) ----------
\@ifundefined{estrelles}{}{%
  \renewcommand{\estrelles}[1]{{\color{curs}\ifcase#1\or\faStar\faStar[regular]\faStar[regular]\or\faStar\faStar\faStar[regular]\or\faStar\faStar\faStar\fi}}}

% ---------- capçalera i peu de les pàgines ----------
% primera pàgina: logos institucionals al peu
\fancypagestyle{primera}{\fancyhf{}\renewcommand{\headrulewidth}{0pt}
  \fancyfoot[L]{\raisebox{-0.32\height}{\includegraphics[height=0.78cm]{logo-gva.png}}\hspace{5mm}%
    \raisebox{-0.32\height}{\includegraphics[height=0.78cm]{logo-vermellar.png}}}
  \fancyfoot[R]{\scriptsize\color{muted}\href{https://creativecommons.org/licenses/by-nc-sa/4.0/deed.ca}{\ccbyncsa}\ \ CC BY-NC-SA 4.0 · fentmates.org\qquad\small\thepage}}
% la resta: el símbol a l'esquerra i el curs a la dreta, amb el títol al centre com abans
\fancyhead[L]{\raisebox{-0.3\height}{\includegraphics[height=0.62cm]{\fm@marca}}}
\fancyhead[R]{\small\bfseries\color{cursfosc}\fm@cursnom}
\renewcommand{\headrule}{\color{curs!45}\hrule height \headrulewidth}
\makeatother

%   \fmcurs{1eso}
% i, al començament del document:
%   \capcalera{Matemàtiques}{1r ESO}{Fitxa d'exercicis · Unitat 2}{Divisibilitat}
% Opcional abans de \input: \def\fmneed{14} (espai mínim abans de cada secció).
% =====================================================================
\makeatletter
\def\fm@logodir{../../logo/pdf}
\providecommand{\fmneed}{10}
\def\fm@cursnom{}

% clau del curs, color i color del text damunt del color
\newcommand{\fm@curs}[3]{\@namedef{fm@color@#1}{#2}\@namedef{fm@text@#1}{#3}}
\fm@curs{1eso}{E6B84A}{navy}
\fm@curs{2eso}{3B82F6}{white}
\fm@curs{3eso}{FF7A50}{white}
\fm@curs{4eso-a}{A78BFA}{white}
\fm@curs{4eso-b}{6D28D9}{white}
\fm@curs{1batx-mi}{9B1C1C}{white}
\fm@curs{1batx-ccss}{EF6B6B}{white}
\fm@curs{2batx-mii}{1E6B47}{white}
\fm@curs{2batx-ccss}{2BB3A3}{white}

\newcommand{\fmcurs}[1]{%
  \@ifundefined{fm@color@#1}{\PackageError{fentmates}{Curs desconegut: #1}{}}{}%
  \edef\fm@tmp{\noexpand\definecolor{curs}{HTML}{\@nameuse{fm@color@#1}}%
               \noexpand\colorlet{curstext}{\@nameuse{fm@text@#1}}}\fm@tmp
  \colorlet{cursfosc}{curs!72!black}%
  \def\fm@logo{\fm@logodir/vertical-#1.pdf}%
  \def\fm@marca{\fm@logodir/marca-#1.pdf}}

% ---------- capçalera de la primera pàgina ----------
% #1 matèria (petit), #2 curs, #3 tipus de document i unitat, #4 títol
\newcommand{\capcalera}[4]{%
  \gdef\fm@cursnom{#2}%
  \par\noindent
  \begin{tcolorbox}[colback=curs,colframe=curs,boxrule=0pt,arc=2.2mm,width=3.5cm,height=2.15cm,
      valign=center,left=3mm,right=2mm,top=1mm,bottom=1mm,nobeforeafter,box align=center]
    \color{curstext}{\footnotesize #1}\par\vspace{1pt}{\fontsize{19}{21}\selectfont\bfseries #2\par}%
  \end{tcolorbox}\hspace{5mm}%
  \begin{minipage}[c]{\dimexpr\linewidth-3.5cm-5mm-2.9cm\relax}
    {\color{cursfosc}\small\bfseries\MakeUppercase{#3}}\par\vspace{2pt}
    {\color{navy}\fontsize{20}{23}\selectfont\bfseries #4\par}
  \end{minipage}\hfill
  \begin{minipage}[c]{2.6cm}\raggedleft\includegraphics[width=2.6cm]{\fm@logo}\end{minipage}\par
  \vspace{2pt}}

% ---------- barres de secció del color del curs ----------
\renewcommand{\seccio}[2]{\par\needspace{\fmneed\baselineskip}\bigskip
  \begin{tcolorbox}[colback=curs,colframe=curs,coltext=curstext,arc=2mm,boxrule=0pt,left=3mm,right=3mm,top=1.2mm,bottom=1.2mm]
  {\large\bfseries #1\quad #2}
  \end{tcolorbox}\vspace{-2pt}}

% ---------- estrelles de dificultat del color del curs (a Batxillerat no se'n posen) ----------
\@ifundefined{estrelles}{}{%
  \renewcommand{\estrelles}[1]{{\color{curs}\ifcase#1\or\faStar\faStar[regular]\faStar[regular]\or\faStar\faStar\faStar[regular]\or\faStar\faStar\faStar\fi}}}

% ---------- capçalera i peu de les pàgines ----------
% primera pàgina: logos institucionals al peu
\fancypagestyle{primera}{\fancyhf{}\renewcommand{\headrulewidth}{0pt}
  \fancyfoot[L]{\raisebox{-0.32\height}{\includegraphics[height=0.78cm]{logo-gva.png}}\hspace{5mm}%
    \raisebox{-0.32\height}{\includegraphics[height=0.78cm]{logo-vermellar.png}}}
  \fancyfoot[R]{\scriptsize\color{muted}\href{https://creativecommons.org/licenses/by-nc-sa/4.0/deed.ca}{\ccbyncsa}\ \ CC BY-NC-SA 4.0 · fentmates.org\qquad\small\thepage}}
% la resta: el símbol a l'esquerra i el curs a la dreta, amb el títol al centre com abans
\fancyhead[L]{\raisebox{-0.3\height}{\includegraphics[height=0.62cm]{\fm@marca}}}
\fancyhead[R]{\small\bfseries\color{cursfosc}\fm@cursnom}
\renewcommand{\headrule}{\color{curs!45}\hrule height \headrulewidth}
\makeatother

\fmcurs{4eso-b}

\begin{document}
\thispagestyle{primera}

\capcalera{Matemáticas B}{4\textsuperscript{o} ESO}{Ficha de ejercicios · Unidad 3}{\raggedright Ecuaciones e inecuaciones}
\vspace{10pt}
Nombre y apellidos: \rule{6.6cm}{0.4pt}\hfill Grupo: \rule{1.6cm}{0.4pt}\hfill Fecha: \rule{2.2cm}{0.4pt}
\vspace{8pt}

\begin{tcolorbox}[colback=soft,colframe=soft,arc=2mm,left=3mm,right=3mm,top=1.5mm,bottom=1.5mm]
\small Ejercicios graduados por niveles de dificultad:\quad \estrelles{1}\ básico \qquad \estrelles{2}\ medio \qquad \estrelles{3}\ avanzado.\\[2pt] Resuélvelos en el cuaderno mostrando todos los pasos. Comprueba siempre las soluciones de las ecuaciones con fracciones, con radicales y logarítmicas, y da las soluciones de las inecuaciones como intervalo y representadas en la recta.
\end{tcolorbox}

% =====================================================
\seccio{1}{Ecuaciones de segundo grado}

\exf{1} Resuelve las ecuaciones incompletas:\begin{tasks}(3)
  \task $x^2-16=0$
  \task $3x^2+12x=0$
  \task $2x^2-50=0$
  \task $x^2+9=0$
  \task $5x^2=3x$
  \task $4x^2-1=0$
\end{tasks}

\ex{1} Resuelve las ecuaciones:\begin{tasks}(3)
  \task $x^2-5x+6=0$
  \task $x^2+2x-15=0$
  \task $2x^2-7x+3=0$
  \task $x^2-6x+9=0$
  \task $x^2+x+3=0$
  \task $6x^2-x-1=0$
\end{tasks}\recorda{$ax^2+bx+c=0\Rightarrow x=\dfrac{-b\pm\sqrt{b^2-4ac}}{2a}$.}

\ex{2} Opera hasta obtener la forma $ax^2+bx+c=0$ y resuelve:\begin{tasks}[after-item-skip=10pt](2)
  \task $(x+2)^2=3x+10$
  \task $x(x-3)=2(x+7)$
  \task $(2x-1)(x+3)-x^2=3$
  \task $\dfrac{x^2}{3}-\dfrac x2=\dfrac32$
\end{tasks}

\ex{3} Elimina los paréntesis y los denominadores, pasa a la forma $ax^2+bx+c=0$ y resuelve:\begin{tasks}[after-item-skip=14pt](2)
  \task $\dfrac{x(x+3)}{2}-\dfrac{(x-1)(x+1)}{3}=4$
  \task $\dfrac{(x+3)(x-3)}{4}-\dfrac{(x-1)^2}{2}=-2$
  \task $\dfrac{(2x+1)(x-3)}{5}-\dfrac{x-1}{2}=\dfrac{x^2-13}{10}$
  \task $\dfrac{(x+1)^2}{2}-\dfrac{(x-1)^2}{3}=\dfrac{x+11}{6}+x$
  \task $\dfrac{x(x-1)}{3}-\dfrac{(x+1)(x-2)}{4}=1$
  \task $\dfrac{(x-1)(x+2)}{4}-\dfrac{x+1}{2}=\dfrac{x-7}{4}$
  \task $\dfrac{(x+2)^2}{4}-\dfrac{x-1}{2}=\dfrac{x^2+11}{8}$
  \task $2(x-1)^2-3(x+1)(x-2)=x(x-5)-4$
\end{tasks}\recorda{Primero desarrolla los productos y las identidades notables; después multiplica todos los términos por el m.c.m. de los denominadores. ¡Cuidado con el signo menos delante de una fracción!}

\ex{2} Sin resolverlas, indica cuántas soluciones tiene cada ecuación:\begin{tasks}(3)
  \task $3x^2-2x+5=0$
  \task $4x^2-12x+9=0$
  \task $x^2+x-1=0$
\end{tasks}Halla después $k$ para que $x^2-6x+k=0$ tenga una solución doble.

\ex{3} Escribe una ecuación de segundo grado que tenga por soluciones $x=3$ y $x=-5$.


% =====================================================
\seccio{2}{Ecuaciones bicuadradas}

\exf{1} Resuelve las ecuaciones bicuadradas:\begin{tasks}(2)
  \task $x^4-10x^2+9=0$
  \task $x^4-5x^2-36=0$
  \task $x^4+3x^2+2=0$
  \task $4x^4-17x^2+4=0$
\end{tasks}\recorda{Haz el cambio $t=x^2$. Cada solución $t>0$ da $x=\pm\sqrt t$; las $t<0$ no dan soluciones.}

\ex{2} Resuelve:\begin{tasks}(3)
  \task $x^4-8x^2=0$
  \task $(x^2-1)(x^2+2)=10$
  \task $x^4-16=0$
\end{tasks}

\ex{3} Opera hasta obtener una ecuación bicuadrada y resuelve:\begin{tasks}[after-item-skip=14pt](2)
  \task $(x^2-2)(x^2+2)=3x^2$
  \task $\dfrac{x^2(x^2-5)}{2}=-2$
  \task $(x^2+1)^2-3(x^2+1)=10$
  \task $\dfrac{x^4}{4}-\dfrac{5x^2}{2}+4=0$
  \task $\dfrac{x^4}{3}-\dfrac{x^2+2}{2}=x^2-\dfrac52$
  \task $(2x^2-1)^2-(x^2+2)^2=3x^4-21$
\end{tasks}


% =====================================================
\seccio{3}{Ecuaciones de grado mayor que 2}

\exf{2} Resuelve con la regla de Ruffini:\begin{tasks}(2)
  \task $x^3-6x^2+11x-6=0$
  \task $x^3+2x^2-x-2=0$
  \task $x^3-3x+2=0$
  \task $2x^3+x^2-8x-4=0$
\end{tasks}\recorda{Prueba con los divisores del término independiente hasta llegar a una ecuación de segundo grado.}

\ex{2} Pasa todos los términos a un miembro, factoriza y resuelve:\begin{tasks}(3)
  \task $x^4-x^3-7x^2+x+6=0$
  \task $x^3=4x$
  \task $x^4+x^3=x^2+x$
\end{tasks}


% =====================================================
\seccio{4}{Ecuaciones factorizadas}

\exf{1} Resuelve sin multiplicar los factores:\begin{tasks}(2)
  \task $(x-4)(x+1)=0$
  \task $x(2x+3)(x-5)=0$
  \task $(x-1)^2(x+7)=0$
  \task $3(x+2)(4x-1)=0$
\end{tasks}\recorda{Un producto es cero si alguno de los factores es cero.}

\ex{2} Resuelve:\begin{tasks}(2)
  \task $(x^2-16)(x+3)=0$
  \task $(x^2-3x+2)(x^2+5)=0$
  \task $x^3(x^2-2x-8)=0$
  \task $(x-2)(x+1)=4$
\end{tasks}


% =====================================================
\seccio{5}{Ecuaciones con fracciones algebraicas}

\exf{1} Resuelve:\begin{tasks}[after-item-skip=12pt](2)
  \task $\dfrac5x-1=4$
  \task $\dfrac{x}{x-1}=\dfrac32$
  \task $\dfrac1x+\dfrac{2}{x+1}=1$
  \task $x+\dfrac4x=5$
\end{tasks}

\ex{2} Resuelve y comprueba las soluciones:\begin{tasks}[after-item-skip=14pt](2)
  \task $\dfrac{x}{x+2}-\dfrac{1}{x-2}=\dfrac{8}{x^2-4}$
  \task $\dfrac{3}{x-1}-\dfrac{2}{x+1}=\dfrac{6}{x^2-1}$
  \task $\dfrac{1}{x-1}+\dfrac{1}{x+1}=\dfrac43$
  \task $\dfrac{x}{x-2}-\dfrac{2}{x+2}=\dfrac{8}{x^2-4}$
  \task $\dfrac{x+3}{x-1}-\dfrac{x-3}{x+1}=4$
  \task $\dfrac2x-\dfrac{3}{x^2+x}=\dfrac{1}{x+1}$
\end{tasks}\recorda{Descarta las soluciones que anulan algún denominador.}

\ex{3} Resuelve $\dfrac{x+1}{x}-\dfrac{x}{x+1}=\dfrac56$.


% =====================================================
\seccio{6}{Ecuaciones con radicales}

\exf{1} Resuelve:\begin{tasks}(4)
  \task $\sqrt{x+5}=3$
  \task $\sqrt{2x-1}=5$
  \task $\sqrt{x^2+9}=5$
  \task $\sqrt{x+3}=-2$
\end{tasks}

\ex{2} Resuelve y comprueba las soluciones:\begin{tasks}(2)
  \task $\sqrt{x+2}=x$
  \task $x+\sqrt{x-1}=7$
  \task $\sqrt{3x+4}-x=0$
  \task $2\sqrt x=x-3$
  \task $\sqrt{2x+3}=x$
  \task $x-\sqrt{x+10}=2$
  \task $\sqrt{5x+4}-2=x$
  \task $3+\sqrt{x-1}=x$
\end{tasks}\recorda{Aísla el radical, eleva al cuadrado y comprueba todas las soluciones en la ecuación inicial.}

\ex{3} Resuelve (hay dos radicales):\begin{tasks}(2)
  \task $\sqrt{x+7}-\sqrt x=1$
  \task $\sqrt{2x+1}+\sqrt{x-3}=4$
\end{tasks}


% =====================================================
\seccio{7}{Ecuaciones exponenciales}

\exf{1} Escribe los dos miembros como potencias de la misma base y resuelve:\begin{tasks}(3)
  \task $2^x=32$
  \task $3^{x-2}=27$
  \task $5^{2x}=\frac{1}{25}$
  \task $4^x=8$
  \task $2^{x^2-1}=8$
  \task $10^{x+1}=0,001$
\end{tasks}

\ex{1} Resuelve tomando logaritmos (da el resultado con tres decimales):\begin{tasks}(4)
  \task $2^x=10$
  \task $3^x=50$
  \task $7^{x-1}=30$
  \task $1,05^x=2$
\end{tasks}\recorda{$a^x=b\Rightarrow x=\dfrac{\log b}{\log a}$.}

\ex{2} Resuelve sacando factor común:\begin{tasks}(3)
  \task $2^{x+1}+2^x=24$
  \task $3^{x+2}-3^x=72$
  \task $5^x+5^{x-1}=30$
\end{tasks}

\ex{3} Resuelve con un cambio de variable:\begin{tasks}(3)
  \task $4^x-5\cdot2^x+4=0$
  \task $9^x-10\cdot3^x+9=0$
  \task $4^x-2^x-12=0$
\end{tasks}


% =====================================================
\seccio{8}{Ecuaciones logarítmicas}

\exf{1} Resuelve aplicando la definición de logaritmo:\begin{tasks}(4)
  \task $\log x=3$
  \task $\log_2(x+1)=4$
  \task $\ln x=0$
  \task $\log_5(3x-1)=2$
\end{tasks}

\ex{2} Resuelve y comprueba las soluciones:\begin{tasks}(2)
  \task $\log x+\log4=\log20$
  \task $\log(x+3)-\log(x-1)=\log3$
  \task $2\log x=\log(5x-6)$
  \task $\log x+\log(x+9)=1$
  \task $\log(x+1)+\log(x-2)=1$
  \task $2\log x-\log(x-2)=\log9$
  \task $\log(2x+4)-\log x=\log3$
  \task $\log_3(x+2)+\log_3(x-4)=3$
\end{tasks}\recorda{Agrupa en un solo logaritmo en cada miembro. Los argumentos de los logaritmos deben ser positivos.}

\ex{3} Resuelve:\begin{tasks}(2)
  \task $\log(x+6)=1+\log(x-3)$
  \task $\log_2x+\log_2(x-2)=3$
\end{tasks}


% =====================================================
\seccio{9}{Inecuaciones de primer grado}

\exf{1} Resuelve y representa la solución:\begin{tasks}(4)
  \task $2x-5>3$
  \task $3-x\le5$
  \task $4x+1<2x-7$
  \task $-3x\ge12$
\end{tasks}\recorda{Si multiplicas o divides por un número negativo, la desigualdad cambia de sentido.}

\ex{2} Resuelve:\begin{tasks}[after-item-skip=12pt](2)
  \task $3(x-2)-2(x+1)\ge x-8$
  \task $\dfrac{x+1}{2}-\dfrac{x-3}{4}<2$
  \task $2x-\dfrac{3-x}{2}>\dfrac x3$
  \task $5x-1>5(x+1)$
\end{tasks}

\ex{2} Con la tarifa A, un móvil cuesta 40 € al mes. Con la tarifa B cuesta 15 € al mes más 0,50 € por cada GB. ¿A partir de cuántos GB es más cara la tarifa B?


% =====================================================
\seccio{10}{Inecuaciones de segundo grado o superior}

\exf{2} Resuelve y representa la solución:\begin{tasks}(3)
  \task $x^2-9<0$
  \task $x^2-3x-10\ge0$
  \task $x^2+4x>0$
  \task $-x^2+2x+3>0$
  \task $x^2-4x+4\le0$
  \task $x^2+x+1>0$
\end{tasks}\recorda{Calcula las raíces, divide la recta en intervalos y estudia el signo en cada uno.}

\ex{3} Resuelve:\begin{tasks}(3)
  \task $x^3-x>0$
  \task $(x-1)(x+2)(x-4)\le0$
  \task $x^3-2x^2-3x\ge0$
\end{tasks}

\ex{2} Los lados de un rectángulo miden $x$ y $x+3$ cm. ¿Para qué valores de $x$ el área es menor que 40 cm$^2$?


% =====================================================
\seccio{11}{Problemas}

\exf{2} El producto de dos números naturales consecutivos es 156. ¿Cuáles son?

\ex{2} Un rectángulo tiene 34 cm de perímetro y 13 cm de diagonal. ¿Cuánto miden los lados?

\ex{2} ¿Cuántos años deben pasar para que 5000 € a un 3\,\% de interés compuesto anual se conviertan en 6000 €?

\ex{2} Un grupo de amigos alquila un apartamento por 600 €. Si hubieran sido dos amigos más, cada uno habría pagado 25 € menos. ¿Cuántos amigos son?

\ex{3} Un grifo tarda 5 horas más que otro en llenar un depósito. Si los abrimos juntos, lo llenan en 6 horas. ¿Cuánto tarda cada uno por separado?

\ex{2} Si a un número positivo le sumamos su raíz cuadrada, obtenemos 42. ¿Cuál es el número?

\ex{2} La suma de un número y su inverso es $\frac{26}{5}$. ¿Cuál es el número?

\ex{3} ¿Para qué números su cuadrado es menor que su triple? ¿Y para cuáles su cubo es mayor que el propio número?

\ex{3} Una población de bacterias sigue la fórmula $N=500\cdot2^{t/3}$, con $t$ en horas. ¿Cuándo habrá 16\,000 bacterias?

\vfill
\begin{center}\color{muted}\small\faCheckCircle[regular]\ \ ¡Ánimo! Comprueba tus resultados con las herramientas interactivas de los apuntes de la unidad.\end{center}

\end{document}
